\documentclass[10pt,a4paper]{amsart}
\usepackage[margin=19mm]{geometry}
\usepackage{amsmath,amssymb,mathtools}
\usepackage[scr=rsfs,cal=euler]{mathalfa}
\usepackage{xfrac}
\usepackage[unicode,hidelinks,bookmarks=false]{hyperref}
\newcommand{\ie}{i.e.\ }
\newcommand{\cf}{cf.\ }
\newcommand{\resp}{respectively\ }
\newcommand{\Subsection}{\subsection}
\providecommand{\nobreakdash}{\nobreak\hbox{-}}
\setlength{\emergencystretch}{2em}
\multlinegap0pt
\usepackage{amssymb}
\usepackage{algorithm}\usepackage{algpseudocode}
\newtheorem{remark}{Remark}
\newtheorem{lemma}{Lemma}
\DeclareMathOperator{\GenRadon}{\mathcal{R}^\omega_{\delta}}
\DeclareMathOperator{\PixelWeight}{\omega^\mathrm{pd}_{\delta_\mathit{s}}}
\DeclareMathOperator{\RR}{\mathbb{R}}
\DeclareMathOperator{\Radon}{\mathcal{R}}
\DeclareMathOperator{\RayRadon}{\mathcal{R}^\mathrm{rd}_{\delta}}
\DeclareMathOperator{\RayWeight}{\omega^\mathrm{rd}_{\delta_\mathit{x}}}
\DeclareMathOperator{\ZZ}{\mathbb{Z}}
\newcommand{\dd}[1]{\,\mathrm{d}{#1}}
\DeclareMathOperator{\imgdom}{\Omega}
\newcommand{\mres}{\mathbin{\vrule height 1.6ex depth 0pt width
0.13ex\vrule height 0.13ex depth 0pt width 1.3ex}}
\DeclareMathOperator{\sinodom}{\mathcal{S}}
\newcommand{\sumphi}{\sum_{q=0}^{N_\phi-1}}
\newcommand{\sums}{\sum_{p=0}^{N_s-1}}
\newcommand{\sumx}{\sum_{i,j=0}^{N_x-1}}
\title{Convergence of ray- and pixel-driven discretization frameworks in the strong operator topology}
\author{Richard Huber}
\date{Source: arXiv:2503.03069v4 (CC BY 4.0)}
\begin{document}
\maketitle

\noindent\emph{Source and licence.} Convergence of ray- and pixel-driven discretization frameworks in the strong operator topology. Version arXiv:2503.03069v4 (CC BY 4.0), distributed by arXiv under the Creative Commons Attribution 4.0 International licence, http://creativecommons.org/licenses/by/4.0/, as recorded in the arXiv licence metadata for this exact version and shown on the abstract page https://arxiv.org/abs/2503.03069v4. Full licence notice: \texttt{licenses/arxiv-2503.03069-CC-BY-4.0.txt}.\par\smallskip

\noindent\textit{Editorial scope.} Counted unit: the article's lemma Lemma\_weight\_length (source0.tex lines 1015--1044). The article's complete original proof of it is delivered with it (source0.tex lines 1049--1088, one \texttt{proof} environment of 40 lines, which the article places away from the statement). No mathematics is written, repaired or re-derived: the statement and the proof are the article's own lines, in the article's own notation, and no retained line is substituted or re-flowed.  Retained uncounted as the source-local context the proof needs: lines 270--273; lines 556--558; lines 677--679; lines 688--690; lines 911--917.  Nothing was omitted from any retained span, so the package carries no omission record.  No retained line contains a citation, so the wrapper carries no bibliography and no bibliography entry.  No retained line contains a figure environment, an image-inclusion command or drawing code, so no image is delivered and the package declares no asset dependency.  Editorial formatting only, in addition: an AMS \texttt{amsart} preamble that restates the article's own declarations and macros used by the retained lines, namely the environments \texttt{remark}, \texttt{lemma} on separate counters, together with the standard packages those lines need.  The retained environments print in this document as Remark 1, Lemma 1 in order of appearance, whereas the article prints the same items with the numbering of its own full document.  The complete native source of the article as distributed in the arXiv e-print for this exact version is bundled unchanged as \texttt{sources/174230/source0.tex}.
\begin{align}
\label{equ_def_radon_transform}
[\Radon f](\phi,s):= \int_{\RR^2} f(x) \dd {\mathcal{H}^{1}\mres L_{\phi,s}}(x) = \int_{\RR} f(s\vartheta_\phi+t \vartheta_\phi^{\perp}) \dd t 
\end{align}

\begin{equation}\label{equ_def_pixel_weight}
\PixelWeight(\phi,t)=\PixelWeight(t):= \frac{1}{\delta_s^2} \max\{\delta_s-|t|,0\} \qquad \text{for } t\in \RR, \  \phi\in \left [0,\pi\right[.
\end{equation}

\begin{equation} \label{equ_lemma_intersectionlength}
\delta_x^2 \RayWeight(\phi,x_{ij}\cdot \vartheta_\phi - s)= \mathcal{H}^1(L_{\phi,s}\cap X_{ij}) - \frac 1 2 \mathcal{H}^1(L_{\phi,s}\cap \partial X_{ij}) .
\end{equation}

    \begin{equation}\label{equ_def_convolutional_Radon_transform}
[\GenRadon f](\phi,s) := \sumphi\sums v_{qp}(\phi,s) \sumx \omega (\phi_q,x_{ij}\cdot \vartheta_q -s_p) \int_{X_{ij}} f(x) \dd x.
    \end{equation}

\begin{remark} \label{remark_using_hausdorff_ae}
Obviously, $u_{ij}=\chi_{X_{ij}}- \frac{1}{2} \chi_{\partial X_{ij}}=\chi_{X_{ij}}$ in a (Lebesgue) almost everywhere sense. Thus, also $\Radon (\chi_{X_{ij}}- \frac{1}{2} \chi_{X_{ij}})= \Radon \chi_{X_{ij}}$ almost everywhere and values of Lebesgue null sets are irrelevant.
However, if we evaluate the Radon transform pointwise for discretization purposes, suddenly Lebesgue null sets in $\imgdom$ and $\sinodom$ could be relevant; we are particularly concerned with the angles in $\frac{\pi}{2}\ZZ$ and $|x_{ij}\cdot \vartheta_q-s_p|=\underline s(\phi)=\delta_x$. Rather, due to the Radon transform's relation to the one-dimensional Hausdorff measure $\mathcal{H}^1$, we define $u_{ij}$ in an $\mathcal{H}^1$ almost everywhere sense, with the value $1$ in the interior of $X_{ij}$, $\frac{1}{2}$ on the boundary. The corners are a bit of a special case as we would actually like the values to be $\frac{1}{4}$, but these are an $\mathcal{H}^1$ null set and thus of no consequence. Analogously, functions $f_\delta\in U_\delta$ are understood as being defined $\mathcal{H}^1$ almost everywhere. Note that $U_\delta$ is not a native subset of $L^2(\imgdom)$, but of $L^2([-1,1]^2)$. We tacitly extend the definition of the Radon transform in \eqref{equ_def_radon_transform} to $L^2([-1,1]^2)$ where necessary. One could make analogous constructions for the $v_{qp}$ basis being defined $\mathcal{H}^1$ almost everywhere, however, these will not be necessary for our analysis below.

%, so our inclusion of the boundary appears meaningless. However, in a discretization we can only evaluate finitely many lines, a nullset that can suddenly very relevant.
%On the contrary, the corners of the $X_{ij}$ are neglected in our considerations (there things will add up to $2$ instead of $1$) since they form a Hausdorff null set and thus do not effect our considerations.
\end{remark}

\begin{lemma}[Exact weights] \label{Lemma_weight_length}
Let $f\in L^2(\imgdom)$, and given $\delta$, let $f_\delta =\sumx f_{ij}u_{ij}\in U_\delta$ with $f_{ij}=\frac{1}{\delta_x^2}\int_{X_{ij}}f \dd x$ (the projection of $f$ onto $U_\delta$). Let $(\phi,s)\in \sinodom$ and let $\hat q\in [N_\phi]$ and $\hat p\in [N_s]$ be such that $(\phi,s)\in \Phi_{\hat q}\times S_{\hat p}$. Then, we have
\begin{equation}
\label{equ_lemma_exact_on_U}
[\RayRadon f](\phi,s) = [\RayRadon f_\delta] (\phi_{\hat q},s_{\hat p})= [\Radon f_\delta] (\phi_{\hat q},s_{\hat p}).  
\tag{$\mathrm{exact}^\mathrm{rd}$}
\end{equation}
(Note that here $f_\delta \in U_\delta$ is understood as defined $\mathcal{H}^1$ almost everywhere (see Remark \ref{remark_using_hausdorff_ae}); thus, the pointwise evaluation $[\Radon f_\delta] (\phi_{\hat q},s_{\hat p})$ is well-defined.)

For fixed $\hat i,\hat j\in [N_x]$, the set $P_{\hat i,\hat j,\hat q}:=\{p\in [N_s] \ \big|\ \PixelWeight(x_{\hat i \hat j}\cdot \vartheta_{\hat q} -s_p)\neq 0\}$ satisfies
\begin{equation} \label{equ_lemma_relevant_p_pd}
P_{\hat i,\hat j,\hat q}  \begin{cases}
=\{\hat p\} & \text{if }x_{\hat i \hat j}\cdot\vartheta_{\hat q}=s_{\hat p}  \text{ for some }\hat p\in [N_s],
\\
=\{\hat p,\hat p+1\} \qquad &\text{if }x_{\hat i \hat j}\cdot\vartheta_{\hat q}\in ]s_{\hat p},s_{\hat p+1}[ \text{ for some }\hat p\in [N_s],
\\
\in \{\{0\},\{N_s-1\},\emptyset\} &  \text{else ( if $x_{\hat i \hat j}\cdot\vartheta_{\hat q} \not \in [s_{0},s_{N_s-1}]$)}.
\end{cases}
\end{equation}
Moreover, we have
\begin{equation}
\label{equ_lemma_pixel_sum_s}
    \sums \PixelWeight(x_{\hat i \hat j}\cdot \vartheta_{\hat q} -s_p) \quad  \begin{cases} =\frac{1}{\delta_s} \qquad & \text{if } x_{\hat i \hat j}\cdot \vartheta_{\hat q}\in [s_0,s_{N_s-1}],
    \\
    \leq \frac{1}{\delta_s} & \text{else}.
\end{cases}    
    \tag{$\mathrm{intpol}^{\mathrm {pd}}$}
\end{equation}

\end{lemma}

\begin{proof}[\textbf{Proof of Lemma \ref{Lemma_weight_length}}]
\underline{\eqref{equ_lemma_exact_on_U}:}  We consider $f_\delta=\sumx f_{ij} u_{ij}\in U_\delta$ with the coefficients $f_{ij}= \frac{1}{\delta_x^2}\int_{X_{ij}}f(x) \dd x$. 
We have $[\RayRadon f] (\phi,s)=[\RayRadon f] (\phi_{\hat q},s_{\hat p})$ for $(\phi,s)\in \Phi_{\hat q}\times S_{\hat p}$ (these  are constant functions on the sinogram pixels). Since $\int_{X_{ij}}f \dd x = \int_{X_{ij}}f_\delta \dd x$, we have  $[\RayRadon f] (\phi,s)=[\RayRadon f_\delta] (\phi_{\hat q},s_{\hat p})$ per definition in \eqref{equ_def_convolutional_Radon_transform}.

Moreover, we calculate
\begin{align*}
[\Radon f_\delta](\phi_{\hat q},s_{\hat p}) & \overset{\text{lin}}{=} \sumx f_{ij} [\Radon u_{ij}](\phi_{\hat q},s_{\hat p})
\\ 
&\underset{\text{def}}{\overset{\text{per}}{=}} \sumx f_{ij} \left ( \int_{\RR^2} \chi_{X_{ij}} \dd {\mathcal{H}^1 \mres L_{\phi_{\hat q},s_{\hat p}}}- \frac 12 \int_{\RR^2} \chi_{\partial X_{ij}} \dd {\mathcal{H}^1 \mres {L_{\phi_{\hat q},s_{\hat p}}}} \right)
\\
&
=\sumx f_{ij}\left ( \mathcal{H}^1(L_{\phi_{\hat q},s_{\hat p}} \cap X_{ij})- \frac{1}{2}\mathcal{H}^1(L_{\phi_{\hat q},s_{\hat p}} \cap \partial X_{ij}) \right)
\\
& \overset{\eqref{equ_lemma_intersectionlength}}{=} \sumx f_{ij} \delta_x^2\RayWeight (\phi_{\hat q},x_{ij}\cdot\vartheta_{\hat q} -s_{\hat p})
 \underset{}{\overset{\text{\eqref{equ_def_convolutional_Radon_transform}}}{=}} [\RayRadon f_\delta](\phi_{\hat q},s_{\hat p}) .
\end{align*}


\underline{ \eqref{equ_lemma_pixel_sum_s}:} Note that $ \PixelWeight(t)\neq 0$ iff  $t\in{]-\delta_s,\delta_s[}$.
If $x_{\hat i \hat j}\cdot\vartheta_{\hat q}=s_{\hat p}$ for some $\hat p \in [N_s]$, then $\PixelWeight(x_{\hat i \hat j}\cdot \vartheta_{\hat q}-s_{\hat p})=\PixelWeight(0)=\frac{1}{\delta_s}$ and 
\begin{equation}
|x_{\hat i \hat j}\cdot\vartheta_{\hat q}-s_{p}|=|s_{\hat p}-s_p| = |\hat p- p|\delta_s\geq \delta_s
\end{equation}
 for $p\neq \hat p$ and therefore $\PixelWeight(x_{\hat i \hat j}\cdot \vartheta_{\hat q}-s_p)=0$, implying \eqref{equ_lemma_relevant_p_pd} and \eqref{equ_lemma_pixel_sum_s}.

If $x_{\hat i \hat j}\cdot\vartheta_{\hat q}\in [s_0,s_{N_s-1}]$, but $x_{\hat i \hat j}\cdot\vartheta_{\hat q}\neq s_p$ for all $p\in [N_s]$, then there is a $\hat p\in [N_s-1]$ with $x_{\hat i \hat j}\cdot\vartheta_{\hat q}\in ]s_{\hat p},s_{\hat p+1}[$. Recall $s_p=s_0+p \delta_s$, and consequently 
\begin{equation} \label{equ_proof_which_p_are_relevant_pd}
|x_{\hat i \hat j}\cdot\vartheta_{\hat q}-s_p|= \min_{p^* \in \{\hat p,\hat p +1\}} |p-p^*|\delta_s+|x_{\hat i \hat j}\cdot\vartheta_{\hat q}-s_{p^*}|\geq \delta_s
\end{equation}
for $p\not \in \{\hat p,\hat p+1\}$ and thus $\PixelWeight(\phi_q,x_{\hat i \hat j}\cdot \vartheta_{\hat q}-s_p)=0$. 
We set $t = x_{\hat i \hat j}\cdot \vartheta_{\hat q}-s_{\hat p+1} \in ]-\delta_s,0[$, and $t+\delta_s = x_{\hat i \hat j}\cdot \vartheta_{\hat q}-s_{\hat p+1} +\delta_s= x_{\hat i \hat j}\cdot \vartheta_{\hat q}-s_{\hat p}$. 
Then, 
\begin{multline}
 \delta_s^2 \sums \PixelWeight(x_{\hat i \hat j}\cdot \vartheta_{\hat q} -s_p)=(\PixelWeight(x_{\hat i \hat j}\cdot \vartheta_{\hat q} -s_{\hat p})+\PixelWeight(x_{\hat i \hat j}\cdot \vartheta_{\hat q} -s_{\hat p+1}))\delta_s^2
\\
 = (\PixelWeight(t)+\PixelWeight(t+\delta_s))\delta_s^2 \underset{\text{}}{\overset{\eqref{equ_def_pixel_weight}}{=}} \delta_s +t + \delta_s - \delta_s-t = \delta_s,
\end{multline} implying \eqref{equ_lemma_pixel_sum_s}.

If $x_{\hat i \hat j}\cdot\vartheta_{\hat q} \not \in [s_0,s_{N_s-1}]$, there is at most one non-zero summand in \eqref{equ_lemma_pixel_sum_s} (via analogous consideration to \eqref{equ_proof_which_p_are_relevant_pd}), which is also bounded by $\frac{1}{\delta_s}$, implying the claim.
\end{proof}

\end{document}
